% Book-local P03.1 completed-algebraic-closure reduction proof,
% 19 September 2026.  No novelty or priority claim is made.

\begin{proof}[\Cref{thm:p031-completed-algebraic-closure-reduction}]
We first prove openness on the ordinary open set $V$.  Fix $y\in V_L$,
put $x=\pi_{L/K}(y)$, and let $N\subseteq V_L$ be an ordinary open
neighborhood of $y$.  Since $K$ is nontrivially valued and $V$ is open in
analytic affine space, there is a strict affinoid domain $X\subseteq V$
such that $x$ lies in the relative interior $\operatorname{Int}(X/V)$.
Compatibility of relative interiors with scalar extension gives
\[
 y\in\pi_{L/K}^{-1}\bigl(\operatorname{Int}(X/V)\bigr)
       \subseteq\operatorname{Int}(X_L/V_L);
\]
see \cite[\S2.11(5)]{GR25Classical}.

The completed-algebraic-closure openness theorem is recorded and applied
in the proof of \cite[Proposition~3.5(9)]{GR25Classical}, where it is
traced to Berkovich's Corollary~1.3.6.  It makes the scalar map
$X_L\to X$ open.  Hence
\[
 \pi_{L/K}\bigl(N\cap\operatorname{Int}(X_L/V_L)\bigr)
       \cap\operatorname{Int}(X/V)
\]
is an ordinary open neighborhood of $x$ in $V$, contained in
$\pi_{L/K}(N)$.  Therefore $\pi_{L/K}:V_L\to V$ is open.

Scalar projection is continuous and surjective by
\cite[\S2.9]{GR25Classical}.  Applying
\Cref{lem:p031-open-pullback-usc} to $a=s$ gives
\[
              u\circ\pi_{L/K}
                 =s^*\circ\pi_{L/K}
                 =(s\circ\pi_{L/K})^*=s_L^*,
\]
which proves \eqref{eq:p031-no-completed-algebraic-closure-defect}.

By \cite[Proposition~4.2(5)]{GR25Classical}, each
$u_i\circ\pi_{L/K}$ is classically psh on $V_L$.  The pulled-back family
remains locally upper bounded and has pointwise supremum $s_L$.
Applying \Cref{thm:p031-affine-ground-field-perron} over $L$ proves that
$u_L=s_L^*$ is $L$-curvewise psh.  Finally, another application of
\cite[Proposition~4.2(5)]{GR25Classical} gives
\[
 u\text{ is classically psh on }V
 \quad\Longleftrightarrow\quad
 u\circ\pi_{L/K}=u_L\text{ is classically psh on }V_L.
\]
This proves all assertions.
\end{proof}
